\section{讨论与结论}
\label{sec:discussion}

\subsection{三个自然推广及其代价}

有限确定性主定理之所以干净，是因为分区集合有限、概率单纯形紧致、纤维条件可以逐点检查。离开这个范围后，交换式仍有意义，但存在性与可检验性需要新条件。

\paperclaim{标准 Borel 空间。}Markov kernel、可测因子和条件分布可以在标准 Borel 空间中处理，相关测度论背景可见\citet{Kallenberg2021}。然而任意等价关系的商未必仍是良好的标准 Borel 空间；最优 kernel 也不再由有限单纯形紧致性自动达到。附录~\ref{app:extensions}只在商可数生成或由到标准 Borel 空间的可测映射给出、并满足所需正则性时陈述因子化结果。

\paperclaim{随机粗粒化。}把确定性 $Q_s$ 替换为 Markov kernel $Q:B\rightsquigarrow Y$ 后，$KQ=QL$ 是标准的随机通道交织式，并与 Markov functions、统计充分性及 Markov category 语言相接\citep{RogersPitman1981,Fritz2020}。但随机 $Q$ 不再对应唯一分区，$L$ 可能不唯一，单时刻通道交换也不自动提供给定耦合下的完整路径过程。

\paperclaim{连续时间。}对有限状态连续时间链，若生成元满足 $GQ_s=Q_s\bar G$，则半群满足 $e^{tG}Q_s=Q_se^{t\bar G}$；反向在 $t=0$ 可微条件下也成立。无限维过程的定义域、强连续性和 martingale problem 则需要额外分析，可参照连续时间 Markov 过程的标准框架\citep{EthierKurtz1986}。本文不把这些解析条件塞回离散主定理。

\subsection{操作意义}

定理~\ref{thm:gamma-characterization} 把“有没有宏观规律”变成两个可分离的问题：首先，候选分组是否远少于操作性微观状态；其次，被放进同一组的状态是否对下一组给出近似相同的预测。只有两者同时成立，表示才既少又自治。故障--修复链证明这种情形可以由具体结构严格推出，并且宏观最小性可以证明；循环反例则说明它绝非有限系统的普遍事实。

对 Transformer，本文支持的是一个条件式结论：若语言确有低复杂度递推预测状态，模型在声明分布上学到了总体下一 token 规律，并且该状态能从完整推理状态读出，那么可构造误差受控的预测状态动力学。它没有证明现实语言必满足这些条件，也没有从输出行为反推出内部参数的低维几何。这个区分使一般判据与神经网络实例能够相接，而不相互替代。

\subsection{开放问题}

后续工作至少包括：为特定模型族寻找可计算而非指数枚举的 $\Gamma_N$ 上界；在有限样本下联合学习近似操作性基线与宏观分区；研究共享宏观核的 $\min_L\sup_K$ 版本；把状态数与描述长度、计算复杂度或统计复杂度结合；以及在增长时间窗上建立收缩或混合条件。

\subsection{结论}

本文的中心结论可以压缩成一句话：

\begin{quote}
宏观规律是否真正“少而自治”，等价于能否找到一种任务保真的分组，使其相对规模和组内预测分歧同时消失；这一点可以在不预设宏观动力学的情况下直接检验。
\end{quote}

有限性保证候选最优值达到，但不保证它非平凡；非平凡性来自具体系统结构。把这两层分开，既保留了可证明的数学结论，也避免把“可定义的宏观表示”误写成“世界必然简单”。
